% concatenated from Numerical_invariants_2025.tex
\documentclass[12pt]{amsart}

%%%%Standard Packages%%%%

\usepackage{amsfonts,amssymb,enumerate,bm,hhline,makecell,array,mathtools}
\usepackage{mathrsfs}
\usepackage{comment}
\usepackage[normalem]{ulem}
\usepackage[colorlinks=true,citecolor=blue, urlcolor=blue, pagebackref, linkcolor=blue]{hyperref}
%option for back references in hyperref: [pagebackref]
\usepackage{tikz-cd}
\usepackage{amscd}
\usepackage[shortlabels]{enumitem}
\setlist[itemize]{noitemsep,nolistsep}
\usepackage{tensor}
\usepackage{tikz}
\usetikzlibrary{matrix,arrows}
\usepackage{extarrows}
\usepackage[toc,page]{appendix}
\usepackage[all,ps,cmtip,rotate]{xy} 
%\usepackage[usenames,dvipsnames]{color}
\usepackage{color}

%Comment before last version
%\usepackage{showkeys}


%%%%PageSettings%%%%

\setlength{\parindent}{1 cm}
\setlength{\textwidth}{16.9 cm}
\setlength{\topmargin} {-1 cm}
\setlength{\evensidemargin}{0 cm}
\setlength{\oddsidemargin}{0 cm}
\setlength{\footskip}{7 mm}
\setlength{\headheight}{2 mm}
\setlength{\textheight}{23.6 cm}
\setlength{\parskip}{2 mm}


%%%%MathbbLetters%%%%

\def\II{\mathbb{I}}
\def\LL{\mathbb{L}}
\def\VV{\mathbb{V}}

\def\DL{\mathbb{L}'}
\def\BVV{\overline{\mathbb{V}}}
\def\BLL{\overline{\mathbb{L}}}
\def\BII{\overline{\mathbb{I}}}


%%%%MathscrLetters%%%%

\def\cI{\mathscr{I}}
\def\cD{\mathscr{D}}
\def\cA{\mathscr{A}}
\def\cB{\mathscr{B}}
\def\cF{\mathscr{F}}
\def\cL{\mathscr{L}}
\def\cO{\mathscr{O}}
\def\cG{\mathscr{G}}
\def\cP{\mathscr{P}}
\def\cH{\mathscr{H}}
\def\cE{\mathscr{E}}
\def\cJ{\mathscr{J}}
\def\cS{\mathscr{S}}
\def\cC{\mathscr{C}}
\def\cM{\mathscr{M}}
\def\cN{\mathscr{N}}
\def\cK{\mathscr{K}}
\def\cT{\mathscr{T}}
\def\cQ{\mathscr{Q}}
\def\cU{\mathscr{U}}
\def\cV{\mathscr{V}}
\def\cW{\mathscr{W}}
\def\cX{\mathscr{X}}
\def\cY{\mathscr{Y}}
\def\cZ{\mathscr{Z}}


%%%%MathbfLetters%%%%

\def\ba{\ensuremath{\mathbf a}}
\def\bb{\ensuremath{\mathbf b}}
\def\ff{\ensuremath{\mathbf f}}
\def\bI{{\bf I}}
\def\bm{\ensuremath{\mathbf m}}
%\def\tt{\ensuremath{\mathbf t}}
\def\bu{\ensuremath{\mathbf u}}
\def\bv{\ensuremath{\mathbf v}}
\def\bbw{\ensuremath{\mathbf w}}
\def\bp{\mathbf p}
\def\bq{\mathbf q}
\def\bpi{\mathbf \pi}
\def\k{\mathbf k}
 
\def\Z{{\bf Z}}
\def\N{{\bf N}}
\def\F{{\bf F}}
\def\C{{\bf C}}
\def\D{{\bf D}}
\def\A{{\bf A}}
\def\B{{\bf B}}
\def\R{{\bf R}}
\def\Q{{\bf Q}}
\def\P{{\bf P}}
\def\PP{{\bf P}}
\def\T{{\bf T}}
\def\bA{\mathbf{A}}
\def\BQ{\mathbf{Q}}
\def\bQ{\mathbf{Q}}
\def\BX{\mathbf{X}}
\def\bX{\mathbf{X}}
\def\BG{\mathbf{G}}
\def\bG{\mathbf{G}}
\def\X{\mathbf{X}}
\def\Y{\mathbf{Y}}

\def\tk{{\widetilde{\mathbf k}}}
\def\Yd{\mathbf{\check{Y}}}
\def\Yt{\mathbf{\hat{Y}}}
 

%%%%MathfrakLetters%%%%

\def\gp{\mathfrak p}
\def\gn{\mathfrak n}
\def\gr{\mathfrak r}


%%%%GreekLetters%%%%

\def\llambda{\ensuremath{\boldsymbol{\lambda}}}
\def\phi{\varphi}


%%%%OtherLetters%%%%

\def\bk{\bar\k}
\def\bmu{\bar\mu}
\def\cYd{\mathscr{\check{Y}}}
\def\cYt{\mathscr{\hat{Y}}}


%%%%MathSymbols%%%%

\DeclareMathOperator{\Alb}{Alb}
\DeclareMathOperator{\Amp}{Amp}
\DeclareMathOperator{\Aut}{Aut}
\DeclareMathOperator{\Bir}{Bir}
\DeclareMathOperator{\Bl}{Bl}
\DeclareMathOperator{\Br}{Br}
\DeclareMathOperator{\ch}{ch}
\DeclareMathOperator{\CH}{CH}
\DeclareMathOperator{\cl}{cl}
\DeclareMathOperator{\Cl}{Cl}
\DeclareMathOperator{\CG}{\mathsf{CG}}
\DeclareMathOperator{\CKGr}{\mathsf{C}_K\mathsf{Gr}}
\DeclareMathOperator{\coh}{coh}
\DeclareMathOperator{\cod}{cod}
\DeclareMathOperator{\codim}{codim}
\DeclareMathOperator{\Cone}{Cone}
\DeclareMathOperator{\Coker}{Coker}
\DeclareMathOperator{\coker}{Coker}
\DeclareMathOperator{\corank}{corank}
\DeclareMathOperator{\Def}{Def}
\DeclareMathOperator{\Diag}{Diag}
\DeclareMathOperator{\Dis}{{Disc}}
\DeclareMathOperator{\Db}{D\textsuperscript{\rm b}}
\DeclareMathOperator{\disc}{disc}
\def\div{\mathop{\rm div}\nolimits}
\DeclareMathOperator{\Eff}{Eff}
\DeclareMathOperator{\End}{End}
\DeclareMathOperator{\epw}{epw}
\DeclareMathOperator{\eval}{ev}
\DeclareMathOperator{\Exc}{Exc}
\DeclareMathOperator{\Ext}{Ext}
\DeclareMathOperator{\Fix}{Fix}
\DeclareMathOperator{\Fl}{\mathsf{Fl}}
\DeclareMathOperator{\GL}{GL}
\DeclareMathOperator{\Gr}{Gr}
\DeclareMathOperator{\Hdg}{Hdg}
\DeclareMathOperator{\HH}{HH}
\DeclareMathOperator{\Hilb}{Hilb}
\DeclareMathOperator{\Hom}{Hom}
\DeclareMathOperator{\Id}{Id}
\DeclareMathOperator{\Int}{Int}
\def\Im{\mathop{\rm Im}\nolimits}
\DeclareMathOperator{\Ind}{Ind}
\DeclareMathOperator{\Ker}{Ker}
\DeclareMathOperator{\Kum}{Kum}
\DeclareMathOperator{\KKK}{K3}
\DeclareMathOperator{\jb}{\ov{J}}
\DeclareMathOperator{\len}{length}
\DeclareMathOperator{\LG}{LG}
\DeclareMathOperator{\LGr}{\mathsf{LGr}}
\DeclareMathOperator{\lin}{\underset{\mathrm lin}{\equiv}}
\DeclareMathOperator{\Mov}{Mov}
\DeclareMathOperator{\mult}{mult}
\DeclareMathOperator{\Nef}{Nef}
\DeclareMathOperator{\Nm}{Nm}
\DeclareMathOperator{\NS}{NS}
\DeclareMathOperator{\num}{\underset{\mathrm num}{\equiv}}
\DeclareMathOperator{\notnum}{\underset{\mathrm num}{\not\equiv}}
\DeclareMathOperator{\OGr}{\mathsf{OGr}}

%\DeclareMathOperator{\otp}{\otimes\!}
\def\otp{}

\DeclareMathOperator{\Pf}{Pf}
\DeclareMathOperator{\PGL}{PGL}
\DeclareMathOperator{\Pic}{Pic}
\DeclareMathOperator{\Pos}{Pos}
\DeclareMathOperator{\Psef}{Psef}
\DeclareMathOperator{\Proj}{Proj}
\def\Re{\mathop{\rm Re}\nolimits}
\def\RlHom{\mathop{\mathbf{R}\mathcal Hom}\nolimits}
\DeclareMathOperator{\rk}{rk}
\DeclareMathOperator{\SO}{SO}
\DeclareMathOperator{\ssp}{sp}
\DeclareMathOperator{\Sp}{Sp}
\DeclareMathOperator{\Supp}{Supp}
\DeclareMathOperator{\suppdet}{Supp_{det}}
\DeclareMathOperator{\Spec}{Spec}
\DeclareMathOperator{\Stab}{Stab}
\DeclareMathOperator{\rank}{rank}
\DeclareMathOperator{\sing}{sing}
\DeclareMathOperator{\Sing}{Sing}
\DeclareMathOperator{\Span}{Span}
\DeclareMathOperator{\Sym}{Sym}
\DeclareMathOperator{\SL}{SL}
\DeclareMathOperator{\td}{td}
\DeclareMathOperator{\tp}{tp}


%%%%ShortNames%%%%

\def\eps{\varepsilon}
\def\ie{\hbox{i.e.}}
\def\eg{\hbox{e.g.}}

\def\vide{\varnothing}
\def\emptyset{\varnothing}

\def\symv{\VV}

\def\ot{\otimes}
\def\op{\oplus}

\def\setminus{\smallsetminus}

\def\isom{\simeq}
\def\cong{\isom}


%%%%Arrows%%%%

\def\lra{\longrightarrow}
\def\llra{\hbox to 10mm{\rightarrowfill}}
\def\lllra{\hbox to 15mm{\rightarrowfill}}

\def\llla{\hbox to 10mm{\leftarrowfill}}
\def\lllla{\hbox to 15mm{\leftarrowfill}}

\def\dra{\dashrightarrow}
\def\thra{\twoheadrightarrow}
\def\hra{\hookrightarrow}

\def\lhra{\ensuremath{\lhook\joinrel\relbar\joinrel\to}}
\def\lllhra{\ensuremath{\lhook\joinrel\relbar\joinrel\lllra}}

\DeclareMathOperator{\isomto}{\stackrel{{}_{\scriptstyle\sim}}{\to}}
\DeclareMathOperator{\isomlto}{\stackrel{{}_{\scriptstyle\sim}}{\leftarrow}}
\DeclareMathOperator{\isomdra}{\stackrel{{}_{\scriptstyle\sim}}{\dra}}
\DeclareMathOperator{\isomlra}{\stackrel{{}_{\scriptstyle\sim}}{\lra}}
\DeclareMathOperator{\isomllra}{\stackrel{{}_{\scriptstyle\sim}}{\longleftarrow}}


%%%%Names%%%%

\def\HK{Hyper-K\"ahler}
\def\HKm{Hyper-K\"ahler manifold}
\def\hk{hyper-K\"ahler}
\def\hKm{hyper-K\"ahler manifold}
\def\av{abelian variety}
\def\avs{abelian varieties}
\def\ppav{principally polarized abelian variety}
\def\pps{principally polarized abelian surface}
\def\ppavs{principally polarized abelian varieties}
\def\pav{polarized abelian variety}
\def\pavs{polarized abelian varieties} 
\def\algc{algebraically closed}


%%%%Environments%%%%

\newtheorem{lemm}{Lemma}[section]
\newtheorem{theo}[lemm]{Theorem}
\newtheorem{coro}[lemm]{Corollary}
\newtheorem{prop}[lemm]{Proposition}
\newtheorem{claim}{Claim}

\theoremstyle{remark}
\newtheorem{defi}[lemm]{Definition}
\newtheorem{rema}[lemm]{Remark}
\newtheorem{conj}[lemm]{Conjecture}
\newtheorem{note}[lemm]{Note}
\newtheorem{notation}[lemm]{Notation}
\newtheorem{exam}[lemm]{Example}
\newtheorem{ques}[lemm]{Question}
\newtheorem{summ}[lemm]{Summary}


%%%%Various%%%%

%Colors
\newcommand{\Cla}[1]{\textcolor{blue}{C: #1}}
\newcommand{\Da}[1]{\textcolor{green}{D: #1}}
\newcommand{\Em}[1]{\textcolor{magenta}{E: #1}}
\newcommand{\Ol}[1]{\textcolor{cyan}{O: #1}}

\newcommand{\blue}[1]{\textcolor{blue}{#1}}
\newcommand{\red}[1]{\textcolor{red}{#1}}
\newcommand{\green}[1]{\textcolor{green}{#1}}

\newcommand{\rem}[1]{\noindent{\color{blue}{#1}}}

%BlankSpace
\def\blank{\underline{\hphantom{A}}}

%MathSubjClass2020
\makeatletter
\@namedef{subjclassname@2020}{%
  \textup{2020} Mathematics Subject Classification}
\makeatother

%TableOfContents
\newcommand{\changelocaltocdepth}[1]{%
  \addtocontents{toc}{\protect\setcounter{tocdepth}{#1}}%
  \setcounter{tocdepth}{#1}%
}

\makeatletter
\def\@tocline#1#2#3#4#5#6#7{\relax
  \ifnum #1>\c@tocdepth % then omit
  \else
    \par \addpenalty\@secpenalty\addvspace{#2}%
    \begingroup \hyphenpenalty\@M
    \@ifempty{#4}{%
      \@tempdima\csname r@tocindent\number#1\endcsname\relax
    }{%
      \@tempdima#4\relax
    }%
    \parindent\z@ \leftskip#3\relax \advance\leftskip\@tempdima\relax
    \rightskip\@pnumwidth plus4em \parfillskip-\@pnumwidth
    #5\leavevmode\hskip-\@tempdima
      \ifcase #1
       \or\or \hskip 1em \or \hskip 2em \else \hskip 3em \fi%
      #6\nobreak\relax
    \dotfill\hbox to\@pnumwidth{\@tocpagenum{#7}}\par
    \nobreak
    \endgroup
  \fi}
\makeatother

\setcounter{tocdepth}{1}

%LocalMacros
\def\bw#1#2{\textstyle{\bigwedge\hskip-0.9mm^{#1}}\hskip0.2mm{#2}}
\def\sbw#1#2{{\bigwedge\hskip-0.9mm^{#1}}\hskip0.1mm{#2}}

\newcommand{\tsi}{{\tilde{\sigma}}}
\newcommand{\bsi}{\bm{\sigma}}

\def\BA{{\overline{A}}}
\def\BW{{\overline{W}}}

\def\Desc{{\mathsf{Desc}}}
\def\PH{{\mathsf{PH}}}

\def\id{\mathsf{id}}
\def\opp{\mathrm{opp}}
\def\ord{\mathrm{ord}}
\def\spe{\mathrm{spe}}

\newcommand{\hY}{\widehat{Y}}
\newcommand{\tY}{\widetilde{Y}}
\newcommand{\tZ}{\widetilde{Z}}

\def\hh{\ensuremath{\mathsf h}}
\def\lll{\ensuremath{\mathsf l}}
\def\mm{\ensuremath{\mathsf m}}
\def\nn{\ensuremath{\mathsf n}}

\def\l{\ensuremath{\lambda}}
\def\M{\ensuremath{\mathsf M}}
\def\K{\ensuremath{\mathsf K}}
\def\IJ{\ensuremath{\mathsf {IJ}}}
\def\J{\ensuremath{\mathsf J}}

\newcommand{\hA}{{\widehat{A}}}

\makeatletter
\let\@wraptoccontribs\wraptoccontribs
\makeatother

 
 
%%%%%%%%%%%%%%%%%%%%%

\title{Numerical invariants of hyper-K\"ahler manifolds}

\begin{document}

\author[O.~Debarre]{Olivier Debarre}
\address{\parbox{0.9\textwidth}{Universit\'e   Paris Cit\'e, CNRS, Institut de Math\'ematiques de Jussieu-Paris rive gauche,\\[1pt]
8 Place Aur\'elie Nemours, 75013 Paris, France
\vspace{1mm}}}
\email{{olivier.debarre@imj-prg.fr}}

 
\address{\parbox{0.9\textwidth}{Shanghai Center for Mathematical Sciences \& School of Mathematical Sciences,\\[1pt] Fudan University, Shanghai 200438, China
\vspace{1mm}}}
\email{chenjiang@fudan.edu.cn}




 
\contrib[with an appendix by]{Chen Jiang}

 

%\subjclass[2020]{\red{..., ...}}
%\keywords{\red{...}}
\thanks{This project has received funding from the European
Research Council (ERC) under the European
Union's Horizon 2020 research and innovation
programme (Project HyperK --- grant agreement 854361).}


\begin{abstract}
We study  various constraints on the Beauville quadratic form and the  Huybrechts--Riemann--Roch  polynomial for \hKm s, mostly in dimension 6 and in the presence of an isotropic class. 

In an appendix, Chen Jiang proves that in general, the  Huybrechts--Riemann--Roch polynomial can  always be written as a linear combination with nonnegative coefficients of certain explicit polynomials with positive coefficients. This implies that the Huybrechts--Riemann--Roch polynomial satisfies a curious   symmetry property.
\end{abstract}

\maketitle

% \tableofcontents
 \setcounter{tocdepth}{1}


%%%%%%%%%%%%%%%%%%%%%

\section{Introduction}

 
A {\em \hKm} is a simply connected compact K\"ahler manifold $X$  whose space of holomorphic 2-forms is spanned by a symplectic form.\ Its dimension  is necessarily an even number $2n$.\ A fundamental tool in the study of \hKm s is the {\em Beauville} form, 
 a canonical integral nondivisible nondegenerate quadratic form $q_X$ on the free abelian group $H^2(X,\Z)$ (\cite[th.~5]{bea}).\ Its signature is $(3,b_2(X)-3)$
 %, it is proportional to the quadratic form
%$$x\longmapsto \int_X \td(X)^{1/2}\, x^2,$$
 and  there is a positive rational number $c_X$ (the {\em Fujiki constant}) such that (\cite[Theorem~4.7]{fuj})
\begin{equation}\label{defqx}
\forall \alpha\in H^2(X,\Z)\qquad \int_X \alpha^{2n}=c_Xq_X(\alpha)^n.
\end{equation}
%  (in dimension~2, the quadratic form~$q_X$ is just the cup-product).\ 
%Moreover, one has $q_X(x)>0$ for all K\"ahler (for example, ample) classes $x\in H^2(X,\Z)$.

There exists a polynomial $ P_{RR,X}(T)$ (the {\em Huybrechts--Riemann--Roch polynomial})  with rational coefficients, leading term $\frac{c_X}{(2n)!}T^n$ and constant term $n+1$, such that, for every line bundle~$L$ on $X$,  one has (\cite[Corollary~3.18]{huy})
\begin{equation}\label{hrr}
\chi(X,L)=P_{RR,X}(q_X(c_1(L)))
.
\end{equation}
The  objects $q_X$, $c_X$, and $ P_{RR,X}(T)$ only depend on  the topology of $X$ and are in particular  deformation invariant.\
 
 In this note, we first prove in Section~\ref{sec:intro} a curious symmetry property for the polynomial $ P_{RR,X}(T)$ (Proposition~\ref{prop15}).\ This property also follows from a strengthening of \cite[Theorem~1.1]{jia} (which says that the polynomial $ P_{RR,X}(T)$ has positive coefficients)
 proved in the appendix by Chen Jiang.
 
 
 
 
  We then study a conjecture made in~\cite[Conjecture~1.4]{dhmv} (and proved in \cite[Theorem~1.5]{dhmv} when~$n=2$) about the possible values of $ P_{RR,X}(T)$ when the quadratic form~$q_X$ represents $0$.\ There exists then a nonzero class $\lll\in H^2(X,\Z)$   such that
$
\int_X\lll^{2n}=0$ and, for any $\mm\in H^2(X,\Z)$, if one writes $\int_X\lll^n\mm^n=a n!$, the number $a$ is necessarily an integer (\cite[Lemma~2.2]{dhmv}).\ The   conjecture deals with the case  $a=1$ (which happens for   $n$th punctual Hilbert schemes of  K3 surfaces and \hKm s of OG10 deformation type).


\begin{conj}[Debarre--Huybrechts--Macr\`i--Voisin]\label{conjnum}
Let $X$ be a \hk\ manifold of dimension $2n$ with  classes $\lll,\mm\in H^2(X,\Z)$ such that
\[
\int_X\lll^{2n}=0\quad\textnormal{and}\quad\int_X\lll^n\mm^n=n!.
\]
Then $c_X=(2n-1)!!$ and the Huybrechts--Riemann--Roch polynomial of $X$  is
\begin{equation}\label{valp}
 P_{RR,X}(T)= \binom{\frac12 T+1+n}{n}.
\end{equation}
\end{conj}


%\begin{rema}
% By \cite[Example~2.12]{beso}, knowing that the polynomial $ P_{RR,X}(T)$  is given by~\eqref{valp}  implies the bound $b_2(X)\le n+17+\frac{12}{n+1}$ on the second Betti number.
%\end{rema}


 Our main result is the following result (Proposition~\ref{prop43}) which almost  proves the conjecture (one would need to additionally prove that the case $n_X=2$ does not happen) in dimension~6 ($n=3$).

\begin{prop}
Let $X$ be a \hk\ manifold of dimension $6$ with  classes $\lll,\mm\in H^2(X,\Z)$ such that
$
\int_X\lll^{6}=0$ and $\int_X\lll^3\mm^3=3!
$.\ We have $q_X(\lll,\mm)=1$, 
%$c_X=15$, $n_X \in\{2,6 \}$, and
  the quadratic form~$q_X$ is even, the Fujiki constant  $c_X$ is $15$, and
$$P_{RR,X}(T)=
  \binom{\frac{T}{2}+4}{3}-\frac{6-n_X}{16}T^2,$$
  where $n_X \in\{2,6 \}$.
  \end{prop}
 
 
 One may make the following more ambitious  conjecture for small positive values of $a$ (it is verified for all known examples of \hk\ manifolds and proved in general when  $n=2$ in~\cite[Theorem~9.3 and Theorem~1.5]{dhmv}).
 
 \begin{conj}\label{c2}
Let $X$ be a \hk\ manifold of dimension $2n$ with  classes $\lll,\mm\in H^2(X,\Z)$ such that
\[
\int_X\lll^{2n}=0\quad\textnormal{and}\quad\int_X\lll^n\mm^n=a n!\ , \quad\textnormal{with }\ a\in\{1,\dots,n\}.
\]
Then $a=1$ and $X$   is   of K3$^{[n]}$ or   OG10 deformation type.
\end{conj}

%Note that the number $a$ is necessarily an integer by \cite[Lemma~2.2]{dhmv}.

Again when $n=3$, we get in Proposition~\ref{prop45} a much weaker result in the case $a=2$ (which, according to Conjecture~\ref{c2}, should not occur at all).

\begin{prop} 
Let $X$ be a \hk\ manifold of dimension $6$ with  classes $\lll,\mm\in H^2(X,\Z)$ such that
$
\int_X\lll^{6}=0$ and $\int_X\lll^3\mm^3=2\cdot 3!
$.\ We have $q_X(\lll,\mm)=1$, 
%$c_X=15$, $n_X \in\{2,6 \}$, and
  the quadratic form~$q_X$ is even,  the Fujiki constant  $c_X$ is $30$, and
  $$P_{RR,X}(T)=
   \frac{1}{24}T^3+\frac{n_X}{8}T^2+
 \bigl(\frac{4}{n_X}+\frac{n_X^2}{12}\bigr)T
 % \frac{13}6T
  +4,$$
  where $n_X \in\{1,2,3,4 \}$.
\end{prop}




 \section{A symmetry property for the Huybrechts--Riemann--Roch polynomial }\label{sec:intro}

 

Let $X$ be a \hk\ manifold of dimension $2n$.\  In \cite[Definition~17]{nw} (see also  \cite[Definition~2.2]{jia}), Nieper-Wi\ss kirchen defined another quadratic form $\lambda_X$ on~$H^2(X,\R)$ (which is {\em not} integral on~$H^2(X,\Z)$).\ It satisfies (see~\cite[(5.18)]{nw})
\begin{equation}\label{fuj2}
\forall \alpha\in H^2(X,\Z)\qquad \frac{1}{(2n)!} \int_X \alpha^{2n}=A_X \lambda_X(\alpha)^n,
\end{equation}
where $A_X\coloneqq\int_X\td^{1/2}(X)$.\ 
% (\cite[(5.12)]{nw}).\ 
By \cite[Proposition~10]{nw} and \cite[Proposition~2.3]{jia}, one can write
\begin{equation*}%\label{defm}
q_X =m_X  \lambda_X,
\end{equation*}
where $m_X$ is a positive rational number, so that (compare~\eqref{defqx} and~\eqref{fuj2})
\begin{equation}\label{cxax}
c_X=\frac{(2n)!A_X}{m_X^n}. 
\end{equation}
We will also set $n_X\coloneqq 2m_X$.\ When $n>1$, one has (\cite[Section~6]{hs}, \cite[Corollary~5.5]{jia})
\begin{equation}\label{td121}
0< A_X <1.
\end{equation}


%\begin{exam}\label{ex11}
%When $X$ is of K3$^{[n]}$ deformation type, one has $c_X=(2n-1)!!$, $n_X= n+3 $, and $A_X=\frac{(n+3)^n}{n!2^{2n}}$.
%
%When $X$ is of Kum$_n$ deformation type, one has $c_X=(n+1)(2n-1)!!$, $n_X= n+1 $, and $A_X=\frac{(n+1)^{n+1}}{n!2^{2n}}$.
% \end{exam}

%As mentioned above, the leading coefficient of the polynomial $P_{RR,X}(T)$ is $a_n=\frac{c_X}{(2n)!}$.\ By \cite[Lemma~5.7]{jia}, the next coefficient $a_{n-1}$ satisfies
%$$a_{n-1}^n=(2n)^na_n^{n-1}A_X=(2na_n)^n 
%$$
%so that

The Hirzebruch--Riemann--Roch theorem~\eqref{hrr}   takes the form 
\begin{equation}\label{hrrq}
\chi(X,L)=\int_X\td(X)\exp (c_1(L))=Q_{RR,X}(\lambda_X(c_1(L))) ,
\end{equation}
where $Q_{RR,X}(T)=P_{RR,X}(m_XT)$.\ 
% and
%% is a polynomial  with positive rational coefficients such that
%\begin{equation}\label{polq}
%Q_{RR,X}(T) = A_X( T^n+2n T^{n-1} )+ \cdots+n+1.
%\end{equation}
%
%%Since $P_{RR,X}(T)=Q_{RR,X}(\frac1{m_X}T)$, we have
%\begin{equation}\label{jc}
%\begin{aligned} 
%Q_{RR,X}(T) &= A_X( T^n+2n T^{n-1} )+ \cdots+n+1,\\
%P_{RR,X}(T)=Q_{RR,X}(\tfrac1{m_X}T)&=A_X\bigl(\tfrac1{m_X^n}T^n +2n \tfrac1{m_X^{n-1}}T^{n-1}  \bigr)+O(T^{n-2})\\
%&=\frac{c_X}{(2n)!} (T+n_X)^n +O(T^{n-2}).
%\end{aligned}
%\end{equation}
The polynomial $Q_{RR,X}(T)$ was computed in \cite[Theorem~5.2]{nw} in terms of the Chern numbers of $X$.\  The formula is
\begin{equation}\label{nwiss}
Q_{RR,X}(T) = \int_X \exp\Bigl(-\sum_{k=1}^{+\infty}\frac{B_{2k}}{2k}  \ch_{2k}(X)T_{2k}\Bigl(\sqrt{\tfrac14T+1}\Bigr)\Bigr),
\end{equation}
where
\begin{itemize}
\item the $B_{2k}$ are the Bernoulli numbers;
\item the $\ch_{2k}\in H^{2k,2k}(X)$ are the homogeneous components of the Chern character of $X$;
\item the $T_{2k}(Y)$ are the (even) Chebyshev polynomials, defined by $T_{2k}( \cos \theta)=\cos 2k \theta$.
\end{itemize}

\bigskip
 
This formula implies curious symmetry relations for the  polynomials $P_{RR,X}(T)$ and $Q_{RR,X}(T)$ for which we have no geometric explanations.

\begin{prop}\label{prop15}
Let $X$ be a \hk\ manifold of dimension $2n$.\  The polynomial $Q_{RR,X}(T)$  satisfies the symmetry relation 
\begin{equation}\label{symm}
Q_{RR,X}(-T-4)=(-1)^n Q_{RR,X}(T).
\end{equation}
Equivalently, 
\begin{equation}\label{symmp}
P_{RR,X}(-T-2n_X)=(-1)^n P_{RR,X}(T).
\end{equation}
\end{prop}

When $n$ is odd, $-n_X$ is therefore a negative rational root of $P_{RR,X}(T)$.\ In all known examples, it is actually an integer (see 
%Example~\ref{ex11} and 
also Lemma~\ref{le2}).

\begin{proof}
Let $P_k$ be the degree $k$ polynomial such that $P_k(T)=T_{2k}\Bigl(\sqrt{\tfrac14T+1} \Bigr)$.\  Set $\cos \theta\coloneqq \sqrt{\tfrac14T+1}$, so that $T=4(\cos^2\theta-1)=-4\sin^2\theta$.\ We compute 
\begin{multline*}
P_k(-T-4)=T_{2k}\Bigl(\sqrt{-\tfrac14T } \Bigr)=T_{2k}(\sin \theta) =T_{2k}( \cos( \theta-\tfrac{\pi}{2}))=\cos(2k( \theta-\tfrac{\pi}{2}))\\
{}=(-1)^k\cos 2k\theta=(-1)^kT_{2k}(\cos\theta)=(-1)^kT_{2k}\Bigl(\sqrt{\tfrac14T+1} \Bigr)=(-1)^kP_k(T).
\end{multline*}
By~\eqref{nwiss}, the polynomial $Q_{RR,X}(T)$ is a $\Q$-linear combination of polynomials of the type
\begin{equation*} 
P_{k_1}(T)  \cdots P_{k_r}(T)  \int_X    \ch_{2k_1}(X)\cdots  \ch_{2k_r}(X) 
\end{equation*}
for $k_1+\cdots +k_r=n$.\  The proposition therefore follows.
\end{proof}


 \begin{rema}\label{rema22}
 The symmetry relation~\eqref{symmp} implies that the polynomial $P_{RR,X}(T)$ is a linear combination with rational coefficients of the polynomials  $(T+n_X)^{n-2j}$, for $0\le j\le n/2$.\ Since its leading coefficient is $\frac{c_X}{(2n)!}$, we can write
 \begin{equation}\label{jc}
\begin{aligned} 
P_{RR,X}(T)&=\frac{c_X}{(2n)!} (T+n_X)^n +O(T^{n-2})\\
Q_{RR,X}(T)=P_{RR,X}(m_XT) &= A_X( T^n+2n T^{n-1} )+ O(T^{n-2}).
\end{aligned}
\end{equation}
The first two coefficients of $P_{RR,X}$ therefore determine $m_X$, $A_X$, and 
$c_X$ (see also \cite[Lemma~5.7]{jia}).
 
 
 Chen Jiang proves in Appendix~\ref{secapp}  that the polynomial $ Q_{RR,X}(T) $ is     a linear combination with {\em nonnegative} rational coefficients of the  polynomials
 $$Q_k(T)\coloneqq\sum_{j=0}^{k}\binom{k+j+1}{2j+1}T^j%=T^n+2nT^{n-1}+\cdots + n+1
 $$
 for $0\le k\le n$ and $n-k$ even.\ These polynomials satisfy the relation~\eqref{symm},\footnote{One has
 $Q_k\bigl(T+\tfrac1{T}-2\bigr) =\sum_{j=0}^kT^{2j-k}$.\ 
In particular, the polynomials $Q_k(T)$  satisfy~\eqref{symm} (change $T$ into~$-T$) and the roots of $Q_k(T)$ are the $k$ negative real numbers $-4\sin^2\frac{j\pi}{2(k+1)}$ for $1\le j\le k$, so that
$$Q_k(T)=\prod_{1\le j\le k}\Bigl(T+ 4\sin^2\frac{j\pi}{2(k+1)}\Bigr)
.$$
} so this much stronger  result implies Proposition~\ref{prop15}. 
%One has
%   \begin{align*} 
%Q_0(T)&=1,\nonumber\\
%Q_1(T)&=T+2,\nonumber\\
%Q_2(T)&=T^2+4T+3=(T+2)^2-1,\nonumber\\
%Q_3(T)&=T^3+6T^2+10T+4=(T+2)^3-2(T+2),\nonumber\\
%Q_4(T)&=T^4+8T^3+21T^2+20T+5=(T+2)^4-3(T+2)^2+1,\\
%Q_5(T)&=T^5+10T^4+36T^3+56T^2+35T+6.\nonumber
%\end{align*}
%In particular,
%   \begin{align*} 
%Q_{RR,X}(T)&=A_XQ_2(T)+(1-A_X)Q_0(T)&\textnormal{when }n=2,\nonumber\\
%Q_{RR,X}(T)&=A_XQ_3(T)+2(1-A_X)Q_1(T)&\textnormal{when }n=3,\nonumber\\
%Q_{RR,X}(T)&=A_XQ_4(T)+B_XQ_2(T)+(5-5A_X-3B_X)Q_0(T)& \textnormal{when }n=4,\\
%Q_{RR,X}(T)&=A_XQ_5(T)+B_XQ_3(T)+(3-3A_X-2B_X)Q_1(T)& \textnormal{when }n=5.\nonumber
%\end{align*}
% Unfortunately, this does not seem very useful for us.
  \end{rema}
 




 
\begin{coro}\label{coro23}
% One has
%  \begin{equation}\label{jc2}
%P_{RR,X}(T)=\frac{c_X}{(2n)!} (T+n_X)^n +O(T^{n-2})
%\end{equation}
%and
%\begin{itemize}
%\item when $n=2$, 
%\begin{equation}\label{prr2}
%\begin{aligned}
% P_{RR,X}(T)&=\frac{c_X}{24}T^2+\frac{c_Xn_X}{12}T+3,\\
%Q_{RR,X}(T)&=\frac{c_Xn_X^2}{96}(T^2+4T)+3,
%\end{aligned}
%\end{equation}
%\item 
When $n=3$, one has
\begin{equation}\label{prr3} 
%P_{RR,X}(T)=  \frac{c_X}{720}T^3+\frac{c_Xn_X}{240}T^2+\Bigl(\frac{4}{n_X}+\frac{c_Xn_X^2}{360}\Bigr)T+4.
%\begin{aligned}
P_{RR,X}(T)= \frac{c_X}{720}(T+n_X)^3 +\Bigl(\frac{4}{n_X}-\frac{c_X}{720}n_X^2\Bigr) (T+n_X).
%Q_{RR,X}(T)&= \frac{c_Xn_X^3}{5760}(T+2)^3 +\Bigl(2-\frac{c_Xn_X^3}{1440}\Bigr) (T+2).
%\end{aligned}
\end{equation}
%\end{itemize}
\end{coro}

\begin{proof}
%It follows from  the symmetry relation~\eqref{symmp} that the polynomial $P_{RR,X}(T)$ is a linear combination  of the polynomials  $(T+n_X)^3$ and $T+n_X$.\ 
%The relation~\eqref{jc2} then follows from the fact that the leading coefficient of $ P_{RR,X}(T)$ is $ \frac{c_X}{(2n)!} $.\ This immediately gives~\eqref{prr2}.\ For~\eqref{prr3}, w
By Remark~\ref{rema22}, we can write
 $$P_{RR,X}(T)=\frac{c_X}{720}(T+n_X)^3+b(T+n_X), 
$$
where $b$  satisfies
$$\frac{c_X}{720}n_X^3+ bn_X=P_{RR,X}(0)
%=\chi(X,\cO_X)
=4,$$
which gives the desired value for $b$.
%$$b=\frac{4}{n_X}-\frac{c_X}{720}n_X^2
%$$
%and implies the corollary.
\end{proof}

For all known examples of \hKm s $X$ of dimension $2n$, one has
%the polynomial $P_{RR,X}(T)$ is one of the   two polynomials
\begin{equation*}%\label{valpo}
P_{RR,X}(T)= \binom{\frac12 T+1+n}{n}\qquad\textnormal{or}\qquad  
 (n+1)\binom{\frac12 T+n}{n}.
\end{equation*}
The roots of both of these polynomials are negative integers (this was conjectured to hold for all \hKm s in \cite[Conjecture~1.3]{jia}).\ In the next two remarks, we discuss what can be said about the reality of the  roots of the polynomial $P_{RR,X}(T)$ (or, equivalenty, of  $Q_{RR,X}(T)$) in dimensions 4 and 6 (when real, the roots are  negative, since both polynomials have positive coefficients).



 \begin{rema}[Real roots, $n=2$]
\label{en2}
When $n=2$, by~\eqref{jc}, we have
\begin{equation*}%\label{n2}
Q_{RR,X}(T) = A_X(T^2+4 T)+3.
\end{equation*}
Easy computations (\cite[Lemma~4.1]{dhmv}) based on \cite[Main Theorem]{gua} give that
%\[
%A_X = 
%%\frac 18 \Bigl(7-\frac 1 {3\cdot 144} c_4(X)\Bigr)=
% \tfrac 1{8\cdot 144} \bigl(992-  4b_2(X)+b_3(X)\bigr).
%\]
%
%
%one has, from the definitions of $\td(X)$ and $A_X$,
%\begin{align*}
%3&=\chi(X,\cO_X)=\td_4(X)= \frac 1 {720}  (3 c_2^2(X) - c_4(X)  ),\\
%A_X &= \frac 1 {8\cdot 720}  ( 7  c_2^2(X) -   4 c_4(X) ),
%\end{align*}
%hence
%\[
%A_X = \frac 18 \Bigl(7-\frac 1 {3\cdot 144} c_4(X)\Bigr).
%\]
%As noted in \cite[Proof of Theorem~2]{gua}, in dimension 4, the Betti numbers determine the Chern numbers: one has for example
%\[
%c_4(X)=3(4b_2(X)+16-b_3(X)).
%\]
%Using \cite[Main Theorem]{gua}, we obtain that 
\begin{itemize}
\item either $b_2(X) =23$ and  $b_3(X) =0$, 
%and the Hodge numbers of $X$ are those of the Hilbert square of a K3 surface, 
in which case 
%$ c_4(X)= 324$ and
 $A_X=\frac{25}{32}$, 
\item or $b_2(X) \le 8$, in which case 
$ \frac{5}{6}\le A_X\le \frac{131}{144}$.
%$c_4(X)\le 144$ and  
%$\frac{5}{6}\le A_X<1$ (the inequality $A_X<1$ follows from~\eqref{td121} but also from the list of possible Betti numbers found in~\cite{gua}).
\end{itemize}
In particular, the  discriminant 
$
%16A_X^2-12A_X=
4A_X(4A_X-3)$
of the polynomial $Q_{RR,X}(T)$ is positive, hence its roots are real. 
\end{rema}

  \begin{rema}[Real roots, $n=3$]
 When $n=3$, we have by Remark~\ref{rema22}
 $$Q_{RR,X}(T) =(T+2)(A_X(T^2+4T)+2).$$
 The roots of this polynomial  are all real if and only if   the discriminant
 $$16A_X^2-8A_X=8A_X(2A_X-1)$$
  is nonnegative, that is, if and only if $A_X\ge  \frac12$.\ The inequality $A_X>  \frac12$ is equivalent to the inequality (2) in \cite{beso}. It implies an upper   bound on $b_2(X)$.\ If $A_X\le  \frac12$, the class $c_2(X)$ is not in  the image of the morphism $\Sym^2\!H^2(X,\Q)\to H^4(X,\Q)$ (the Verbitsky component).
 
 
 \end{rema}

%\begin{rema}\label{rem22}
%One can express the polynomial $Q_{RR,X}(T) $ in terms of the Chern numbers of $X$ (\cite[(5.12)]{nw}).\ When $n=2$, one has
%$$
%Q_{RR,X}(T)=\frac{1}{5760}(7c_2^2-4c_4)(T^2+4T)+\frac{1}{720}(3c_2^2-c_4).
%$$
%Comparing with Corollary~\ref{coro23}, we obtain
%$$7c_2^2-4c_4= 60 c_Xn_X^2\quad,\quad 3c_2^2-c_4=2160,
%$$
%hence
%$$ c_2^2=12(144- c_Xn_X^2)\quad,\quad  c_4=36(84-c_Xn_X^2)\quad,\quad  \ch_4=5(c_Xn_X^2-72).
%$$
%
%
%
%When $n=3$, we have
%$$
%\begin{aligned}
%Q_{RR,X}(T)= &\frac{1}{30420\cdot 32}(31c_2^3-44c_2c_4+16c_6)(T^3+6T^2)+\frac{1}{30420\cdot 4}(41c_2^3-53 c_2c_4+18c_6)T  \\
%&+ \frac{1}{30420\cdot 2}(10c_2^3-9c_2c_4+2c_6) .
%\end{aligned}
%$$
%Note that $-2$ is always a root of this polynomial.\ 
%Comparing with Corollary~\ref{coro23}, we obtain only two equations
%$$31c_2^3-44c_2c_4+16c_6= 169 c_Xn_X^3\quad,\quad 10c_2^3-9c_2c_4+2c_6=243360.
%$$
%
%\end{rema}




\section{Coefficients of the Huybrechts--Riemann--Roch polynomial}

For each positive integer $n$, we define the positive integer  
\[
C_n\coloneqq\gcd_{r_0,\dots,r_n\in \Z} \prod_{0\le j<k\le n}(r_j^2-r_k^2).
\]
One computes easily $C_1=1$, $C_2=12$, and, with a computer,\footnote{Many thanks to Jieao Song for making these computations.\ For any positive integer $n$, the   primes $p$ that divide~$C_n$ are exactly those such that $p\le  2n-1$ (this is because one can find $n+1$ distinct squares modulo $p$ if and only if $p>2n$).} 
\begin{align*}
C_3&=2^5\cdot 3^3\cdot 5,\\
C_4&= 2^{11} \cdot 3^5 \cdot 5^2 \cdot 7,\\
C_5&= 2^{18} \cdot 3^9 \cdot 5^4 \cdot 7^2,
\\
C_6&= 2^{27} \cdot 3^{14} \cdot 5^6 \cdot 7^3 \cdot 11,
\\
C_7&= 2^{37} \cdot 3^{19} \cdot 5^8 \cdot 7^5\cdot 11^2\cdot 13.
\end{align*}
Let $X$ be a \hk\ manifold of dimension $2n$.\  We write the Huybrechts--Riemann--Roch polynomial as
   \begin{equation*}%\label{jc2n}
 P_{RR,X}(T) =:a_nT^n+\cdots +a_1T+a_0,
 \end{equation*}
 where $a_n= \frac{c_X}{(2n)!}$ and $a_0=n+1$.\ The proof of the following proposition uses the fact that the polynomial $ P_{RR,X}(T)$ takes integral values on every integer represented by $q_X$: this is because of the relation~\eqref{hrr} and the fact that, for every $\alpha\in H^2(X,\Z)$, there is a deformation of $X$ on which $\alpha$ becomes the first Chern class of a line bundle.

\begin{prop}\label{prop12}
Let $X$ be a \hk\ manifold of dimension $2n$.\  
For each $i\in\{0,\dots,n\}$, the coefficient $a_i$ of the polynomial $P_{RR,X}(T)$ belongs to $\frac{1}{2^iC_n}\Z$ (and to $\frac{1}{C_n}\Z$ if the quadratic form~$q_X$ is not even).\ 
In particular, the Fujiki constant $c_X$ is in $ \frac{(2n)!}{2^nC_n}\Z$.
\end{prop}

\begin{proof}
Let $q $ be an integer represented by $q_X$.\ 
For all $r_0,\dots,r_n\in \Z$, the integers $r_0^2q,\dots,r_n^2q$ are also represented by $q_X$, so that $P_{RR,X}(r_j^2q)=\sum_{i=0}^n  a_i r_j^{2i}q^i$ is an integer for all $j\in\{0,\dots,n\}$.\ 
The corresponding linear system with unknowns $a_0q^0,\dots,a_nq^n$ has a   Vandermonde matrix $(r_j^{2i})$, so we get
\[
a_iq^i\prod_{0\le j<k\le n}( r_j^2- r_k^2)   \in\Z
\]
for all $i\in\{0,\dots,n\}$, which implies $a_iq^iC_n\in\Z$.\ 
Since the  integral bilinear form associated with $q_X$ is  not  divisible, the gcd of all integers $q$ represented by $q_X$ is either $2$ (if the form $q_X$ is even) or $1$ (if it is not) and the proposition follows.
\end{proof}

In particular, we get $c_X\in\frac12\Z$ when $n=2$, and $c_X\in\frac1{48}\Z$ when $n=3$.\ For any $n$, Proposition~\ref{prop12} gives the lower bound $c_X\ge \frac{(2n)!}{2^nC_n}$, but what would be really interesting, in order to prove boundedness properties for \hk\ manifolds, would be to find  an upper bound on $c_X$ (see \cite{huyf}).

\begin{rema}\label{rem12}
Assume that $q_X$ represents all large enough even numbers (this is the case for all known examples).\ 
Then $P_{RR,X}(T)$ takes integral values on all large enough even numbers and this implies that its leading coefficient is in $\frac1{n!2^n}\Z$, hence $c_X\in (2n-1)!!\Z$.
\end{rema}


 

 




 




 \section{The Huybrechts--Riemann--Roch polynomial in the presence of an isotropic class}
 
 Let $X$ be a \hk\ manifold of dimension $2n$.\   Assume that there is an isotropic class $\lll\in H^2(X,\Z)$, that is, $q_X(\lll)=0$.\ For any   
  $ \mm\in H^2(X,\Z)$,
 $$a(\mm)\coloneqq\frac{1}{n!}\int_X\lll^n\mm^n$$
 is an integer (\cite[Lemma~2.2]{dhmv}) and 
\begin{equation} \label{cxa}
c_Xq_X(\lll,\mm)^n=a(\mm)\frac{(2n)!}{2^nn!}=a(\mm)(2n-1)!!.
\end{equation}
From now on, we assume $q_X(\lll,\mm)>0$.\ Using~\eqref{cxax} and~\eqref{td121}, we obtain
$$
m_X^n =\frac{(2n)!A_X}{c_X}<\frac{(2n)! }{c_X} = \frac{ 2^nn!q_X(\lll,\mm)^n}{a(\mm) }$$
hence
\begin{equation}\label{mxroot}
m_X < 2q_X(\lll,\mm) \sqrt[\leftroot{-3}\uproot{16}  {\scriptstyle {n}}]{\frac{  n!}{a(\mm) }}.
\end{equation}
 Using the bound $c_X\ge \frac{(2n)!}{2^nC_n}$, we also get $m_X< 2 \sqrt[\leftroot{-1}\uproot{2} n]{C_n}$.

\begin{lemm}\label{le1}
We have
$$
n! q_X(\lll,\mm)^n \mid a(\mm)C_n$$
and,  if $q_X$ is not even,
$$n! 2^n q_X(\lll,\mm)^n \mid a(\mm)C_n.$$

\end{lemm}

%The second statement does not seem to be very useful.

\begin{proof}
Using~\eqref{cxa}, we get
$$
a_n= \frac{c_X}{(2n)!} =\frac{a(\mm)}{2^nn!q_X(\lll,\mm)^n}.
%\quad\textnormal{and}\quad a_{n-1}= \frac{2nm_Xc_X}{(2n)!} =\frac{a(\mm)m_X}{2^{n-1}(n-1)!q_X(\lll,\mm)^n}.
$$
Then use Proposition~\ref{prop12}.\end{proof}

\begin{lemm}\label{le2}
 We have 
\begin{equation*}\label{eq}
a(\mm) \Bigl( \frac{q_X(\mm)+n_X}{2q_X(\lll,\mm)  }-\frac{n-1}2   \Bigr)     \in \Z.
\end{equation*}
In particular,
$$n_X\in   \Z+ \frac{2q_X(\lll,\mm) }{a(\mm)}\Z $$
so that $n_X $ is an integer when $ a(\mm)\in \{1,2\}$.
\end{lemm}

\begin{proof}
For every $t\in\Z$, the number
$$P(t )\coloneqq P_{RR,X}(q_X(t\lll+\mm))=P_{RR,X}(2tq_X(\lll,\mm)+q_X(\mm))
$$
is an integer.\ We have, using~\eqref{jc} and~\eqref{cxa},
   \begin{align*}\label{jc3}
P(t)&=\frac{c_X}{(2n)!} (2tq_X(\lll,\mm)+q_X(\mm)+n_X)^n +O(t^{n-2})\\
&= \frac{a(\mm) }{q_X(\lll,\mm)^n2^nn!} (2^n q_X(\lll,\mm)^nt^n+ n2^{n-1} q_X(\lll,\mm)^{n-1}(q_X(\mm)+n_X)t^{n-1}\bigr)  +O(t^{n-2})\\
&= \frac{a(\mm)}{  n!}    t^n+  \frac{a(\mm)}{q_X(\lll,\mm) 2 (n-1)!}    (q_X(\mm)+n_X)t^{n-1}  +O(t^{n-2}).
\end{align*}
This is an integer for all $t\in\Z$, hence so is
   \begin{align*} 
P(t)-a(\mm)\binom{t+n-1}{n}&=P(t)- a(\mm)\, \frac{ t^n+\frac{n(n-1)}2 t^{n-1}}{  n!}    +O(t^{n-2})\\
&= \Bigl( \frac{q_X(\mm)+n_X}{2q_X(\lll,\mm)  }-\frac{n-1}2   \Bigr)    \frac{a(\mm)}{ (n-1)!}  t^{n-1}   +O(t^{n-2}) .
\end{align*}
This implies the lemma.
\end{proof}

\subsection{Case  $a(\mm)=1$}  We know from~\cite[Theorem~1.5]{dhmv}
%, which uses   methods that rely on Guan's results in \cite{gua}, 
that in dimension 4, this case only  occurs when $X$ is of K3$^{[2]}$ deformation type.\ In particular, $P_{RR,X}(T)$ is then given by~\eqref{valp}.\ We believe (Conjecture~\ref{conjnum}) that the same should happen for any $n\ge2 $ (one would then have   $c_n=(2n-1)!!$ and $n_X=n+3$).\ We study the case $n=3$.


\begin{prop}\label{prop43}
%Let $X$ be a \hKm\ of dimension $6$ with  classes $\lll,\mm\in H^2(X,\Z)$ such that $\int_X\lll^{6}=0$ and $\int_X\lll^3\mm^3=3!$.\  Then, 
Assume $n=3$ and $a(\mm)=1$.\ Then
%\begin{enumerate}[{\rm (a)}]
%\item If $n=2$, we have $q_X(\lll,\mm)=1$, $c_X=3$, 
% $n_X\in \{1,3,5\}$, 
%and the quadratic form $q_X$ is even.\ One also has 
%$$P_{RR,X}(T)=\frac18T^2+\frac{n_X}{4}T+3= \binom{\frac{T}{2}+3}{2}-\frac{5-n_X}{4}T.$$
%\item If $n=3$, we have 
$q_X(\lll,\mm)=1$, $c_X=15$, $n_X \in\{2,6 \}$, and the quadratic form $q_X$ is even.\ One also has 
$$P_{RR,X}(T)=
%\frac18T^2+\frac{n_X}{4}T+3
  \frac{1}{48}T^3+\frac{n_X}{16}T^2+
  %\Bigl(\frac{4}{n_X}+\frac{n_X^2}{24}\Bigr)T
  \frac{13}6T+4
  = \binom{\frac{T}{2}+4}{3}-\frac{6-n_X}{16}T^2$$
%\end{enumerate}
%In both cases, 
and the sublattice $\Z\lll\oplus \Z\mm$ of $(H^2(X,\Z),q_X)$ is  a hyperbolic plane.
\end{prop}

\begin{proof}
% \noindent (a) Assume $n=2$.\ We have $C_2=12$ and we obtain   from Lemma~\ref{le1}
%\begin{equation*}\label{eqdiv}
%q_X(\lll,\mm)^2\mid 6  \qquad\textnormal{(and $2q_X(\lll,\mm)^2 \mid 3 $ if $q_X$ is not even)},
%\end{equation*}
%so that $q_X(\lll,\mm)=1$, $q_X$ is even,  $c_X=3$ from~\eqref{cxa}, and $m_X\in \Z+\frac12$ from Lemma~\ref{le2}.\ We can then use~\eqref{mxroot} to deduce
% $m_X<2\sqrt2
%$  hence the possible values for $n_X=2m_X$.\ The formula for~$P_{RR,X}(T)$ then follows from~\eqref{jc}.
% 
%\bigskip
% \noindent (b) Assume $n=3$.\ 
 We have $C_3=2^5\cdot 3^3\cdot 5$ and we obtain   from Lemma~\ref{le1}
\begin{equation*}\label{eqdiv3}
q_X(\lll,\mm)^3\mid    2^4\cdot 3^2\cdot 5 \qquad\textnormal{(and $ q_X(\lll,\mm)^3 \mid 2 \cdot 3^2\cdot 5 $ if $q_X$ is not even),}
\end{equation*}
so that $ q_X(\lll,\mm)\in\{1,2\}$ (and $ q_X(\lll,\mm)=1 $ if $q_X$ is not even).
 
 \medskip
\noindent{\bf{Assume  $ q_X(\lll,\mm)=1 $.}}\ We have $c_X=15$ from~\eqref{cxa}, Lemma~\ref{le2} gives $q_X(\mm)+n_X\in 2\Z$, and~\eqref{mxroot} gives $m_X<2\sqrt[3]{6}\sim 3.6 $, so that $n_X \in\{1,2,3,4,5,6,7\}$.\ Furthermore, we have, by Corollary~\ref{coro23}, 
$$P_{RR,X}(T)=\frac{1}{48}(T+n_X)^3+\Bigl(\frac{4}{n_X}-\frac{1}{48}n_X^2\Bigr)(T+n_X).
$$
 For all values $q$ taken by $q_X$, this must be an integer when $T=q$, so that
\begin{equation}\label{eqdiv3b}
48n_X\mid   n_X(q+n_X)^3+ (192- n_X^3 )(q+n_X).
\end{equation}
 In particular, $n_X\mid 192 q$.\ If $n_X\in\{5,7\}$, this implies $n_X\mid  q$, which is impossible because   the gcd of all integers $q$ represented by $q_X$ is either $1$ or $2$.\ Otherwise, $16n_X\mid 192$, hence we obtain
\begin{equation}\label{eqdiv3bb}
16 \mid  (q+n_X)^3 - n_X^2  (q+n_X)=q(q+n_X)(q+2n_X).
\end{equation}
\begin{itemize}
\item When $n_X=1$, the relation~\eqref{eqdiv3bb} is equivalent to  $q\equiv 0,6,8,14,15\pmod{16}$.\ The case $q\equiv  15\pmod{16}$ is impossible since $4q$ is also represented but not in this list, hence $q\equiv 0,6,8,14\pmod{16}$ and $q_X$ is even.\ This contradicts the fact that $q_X(\mm)+n_X$ is even.
\item When $n_X=2$,  the relation~\eqref{eqdiv3bb} is equivalent to   $q $ even.
\item When $n_X=3$, the only possible odd value is $q\equiv 13\pmod{16}$.\ This implies that  $4q\equiv 4\pmod{16}$ should also be represented, but $4$ does not satisfy the relation~\eqref{eqdiv3bb}.\ So $q_X$ is even, which contradicts the fact that $q_X(\mm)+n_X$ is even.
\item When $n_X=4$,  the relation~\eqref{eqdiv3bb} is equivalent to  $4\mid q$, which is impossible because   the gcd of all integers $q$ represented by $q_X$ is either $1$ or $2$. 
\item When $n_X=6$,  the relation~\eqref{eqdiv3bb}  is equivalent to  $q $ even.
 \end{itemize}
All in all, we get $n_X \in\{2,6 \}$ and $q_X$ even.
 

\medskip

\noindent{\bf{Assume  $ q_X(\lll,\mm)=2$.}}\ The quadratic form $q_X$ is even, we have $c_X=\frac{15}8$ from~\eqref{cxa},   Lemma~\ref{le2} gives $\frac12 q_X(\mm)+ m_X\in 2\Z$, so that $m_X$ is an integer, and~\eqref{mxroot} gives $m_X<4\sqrt[3]{6}<8$, so that $m_X \in\{1,2,3,4,5,6,7\}$.\ As above, we  deduce from~\eqref{prr3} that
 $$ \frac{1}{8\cdot 48}(2q+ 2m_X)^3+\Bigl(\frac{2}{m_X}-\frac{1}{8\cdot 48}4m_X^2\Bigr)(2q+ 2m_X) $$
is an integer for all values $2q$ taken by $q_X$, so that 
$$48m_X\mid    m_X(q+m_X)^3+ (192- m_X^3 )(q+m_X).
$$
This is ``the same'' relation as~\eqref{eqdiv3b} and the discussion above allows us to conclude that $q$ must be even, so that all values taken by $q_X$ are divisible by $4$.\ This is   impossible because   the gcd of all values taken  by $q_X$ is  $2$.\ So this case does not occur. 
%
%When $n=4$, we have $q_X(\lll,\mm) \mid  12$   (and $ q_X(\lll,\mm)  \mid 6 $ if $q_X$ is not even).\ This leaves many choices.
\end{proof}

%\begin{rema}
%Under the hypotheses of Proposition~\ref{prop43}, we can compute, using Remark~\ref{rem22}, some Chern numbers of $X$.\   When $n=2$, we get
%$$c_2^2=36(48-n_X^2)\ ,\ \ c_4=108(28-n_X^2)\ ,\ \ \ch_4= 15(n_X^2-24) 
%$$
%and when $n=3$,  
%$$31c_2^3-44c_2c_4+16c_6= 1035n_X^3\quad,\quad 10c_2^3-9c_2c_4+2c_6=243360.
%$$
%\end{rema}

%\begin{rema}
%Following \cite{beso}, we  examine whether  $\frac{2na_na_{n-2}}{(n-1)a_{n-1}^2}$ is smaller than $1$.\ When   $n=3$ and  $a(\mm)=1$, this   ratio  is equal to
%$$\frac{6\frac{1}{ 48}\frac{13}{ 6}}{2\frac{n_X^2}{16^2}}= \frac{104}{ 3n_X^2   } $$
%so it smaller than 1 only when $n_X=6$.\ In that case, by \cite[Theorem~1.1]{beso}, we have  
%$$b_2(X)\le 23 
%$$ 
%and there is equality if and only if $c_2(X)\in SH^2(X,\R)$.\ When $n_X=2$, one has  $c_2(X)\notin SH^2(X,\R)$.
% \end{rema}


\subsection{Case  $a(\mm)=2$}   We believe (Conjecture~\ref{c2})  this case should not     occur for any~$n\ge2$ and we know from~\cite[Theorem~9.3]{dhmv} that it does not when $n=2$.\ We study the case $n=3$.


\begin{prop}\label{prop45}
Assume $n=3$ and $a(\mm)=2$.\ Then,
%\begin{enumerate}[{\rm (a)}]
%\item If $n=2$, we have $q_X(\lll,\mm)=1$, $c_X=6$, 
% $n_X\in \{1,2,3\}$, 
%and the quadratic form $q_X$ is even.\ One also has $P_{RR,X}(T)=\frac14T^2+\frac{n_X}{2}T+3$.
%\item If $n=3$, 
 $q_X(\lll,\mm)=1$, $c_X=30$, $n_X \in\{1,2,3,4 \}$, and the quadratic form $q_X$ is even.\ One also has $P_{RR,X}(T)=
   \frac{1}{24}T^3+\frac{n_X}{8}T^2+
 \bigl(\frac{4}{n_X}+\frac{n_X^2}{12}\bigr)T
 % \frac{13}6T
  +4$
%\end{enumerate}
%In both cases, 
and the sublattice $\Z\lll\oplus \Z\mm$ of $(H^2(X,\Z),q_X)$ is a hyperbolic plane.
\end{prop}

\begin{proof}
% \noindent (a) Assume $n=2$.\ We have $C_2=12$ and we obtain   from Lemma~\ref{le1}
%\begin{equation*} 
%q_X(\lll,\mm)^2\mid 12  \qquad\textnormal{(and $ q_X(\lll,\mm)^2 \mid 3 $ if $q_X$ is not even)},
%\end{equation*}
%so that $ q_X(\lll,\mm)\in\{1,2\}$ (and $ q_X(\lll,\mm)=1 $ if $q_X$ is not even).
%
%\medskip
%
%\noindent{\bf{Assume  $ q_X(\lll,\mm)=1 $.}}\ We have $c_X=6$ from~\eqref{cxa}, Lemma~\ref{le2} gives $n_X\in  \Z$, and~\eqref{mxroot} gives $m_X<2 $, so that $n_X \in\{1,2,3\}$.\ Furthermore, we have, by~\eqref{prr2},
%$$P_{RR,X}(T)=\frac14T^2+\frac12n_XT+3, 
%$$
%so that $\frac14 q^2+\frac12n_Xq=\frac14 q(q+2n_X)$ must be an integer for all values $ q$ taken by $q_X$.\ This implies that $q$ must be even.
%
%\medskip
%
%\noindent{\bf{Assume  $ q_X(\lll,\mm)=2 $.}}\ The quadratic form $q_X$ is even, we have $c_X=\frac32$ from~\eqref{cxa}, Lemma~\ref{le2} gives $m_X\in  \Z$, and~\eqref{mxroot} gives $m_X<4 $, so that $m_X \in\{1,2,3\}$.\ As above, $\frac1{16} (2q)^2+\frac18n_X(2q)=\frac14 q(q+2m_X)$ must be an integer for all values $ 2q$ taken by $q_X$.\ This implies that $q$ must be even and contradicts the fact that the gcd of all integers $2q$ represented by $q_X$ is $2$.\ So this case does not happen. 
%
%\bigskip
% \noindent (b) Assume $n=3$.\ 
 We have $C_3=2^5\cdot 3^3\cdot 5$ and we obtain   from Lemma~\ref{le1}
\begin{equation*} 
q_X(\lll,\mm)^3\mid    2^5\cdot 3^2\cdot 5 \qquad\textnormal{(and $ q_X(\lll,\mm)^3 \mid 2^2 \cdot 3^2\cdot 5 $ if $q_X$ is not even),}
\end{equation*}
so that $ q_X(\lll,\mm)\in\{1,2\}$ (and $ q_X(\lll,\mm)=1 $ if $q_X$ is not even).

\medskip

\noindent{\bf{Assume  $ q_X(\lll,\mm)=1 $.}}\ We have $c_X=30$ from~\eqref{cxa}, Lemma~\ref{le2} gives $ n_X\in \Z$, and~\eqref{mxroot} gives $m_X<2\sqrt[3]{3}\sim 2.9 $, so that $n_X \in\{1,2,3,4,5\}$.\ Furthermore, we have, by~\eqref{prr3}, 
$$P_{RR,X}(T)=\frac{1}{24}(T+n_X)^3+\Bigl(\frac{4}{n_X}-\frac{1}{24}n_X^2\Bigr)(T+n_X),
$$
 For all values $q$ taken by $q_X$, this must be an integer when $T=q$, so that
\begin{equation*}%\label{eqdiv3c}
24n_X\mid   n_X(q+n_X)^3+ (96- n_X^3 )(q+n_X).
\end{equation*}
 In particular, $n_X\mid 96 q$.\ If $n_X=5$, this implies $n_X\mid  q$, which is impossible because   the gcd of all integers $q$ represented by $q_X$ is either $1$ or $2$.\ Otherwise, $8n_X\mid 96$, hence we obtain
\begin{equation}\label{eqdiv3cc}
8 \mid  (q+n_X)^3 - n_X^2  (q+n_X)=q(q+n_X)(q+2n_X).
\end{equation}
\begin{itemize}
\item When $n_X=1$, the relation~\eqref{eqdiv3cc} is equivalent to  $q\equiv 0,2,4,6,7\pmod{8}$; this means that every odd value taken by $q_X$ is $\equiv  7\pmod{8}$.\ Assume there exists $\alpha$ such that    $q_X(\alpha)\equiv  7\pmod{8}$.\ Since $q_X(k\lll+\mm)= 2k+q_X(\mm)$, the integer $q_X(\mm)$ must be even and we may even assume $q_X(\mm)=0$.\    For all   $t,u\in\Z$, the integer
$$q_X(t\lll+u\mm+\alpha)=2tu+2tq_X(\lll,\alpha)+2uq_X(\mm,\alpha)+q_X(\alpha)
$$
is odd, hence is $\equiv  7\pmod{8}$.\ This implies $tu+ tq_X(\lll,\alpha)+ uq_X(\mm,\alpha)\equiv 0\pmod4$.\ Taking $t=1$ and $u=0$, we obtain $q_X(\lll,\alpha) \equiv 0\pmod4$; taking $t=0$ and $u=1$, we obtain $q_X(\mm,\alpha) \equiv 0\pmod4$; taking $t=u=1$, we obtain a contradiction.\ Hence $q\equiv 0,2,4,6\pmod{8}$ and $q_X$ is even.\ 
\item When $n_X=2$, the relation~\eqref{eqdiv3cc} is equivalent to  $q\equiv 0,2,4,6\pmod{8}$ and $q_X$ is even.
\item When $n_X=3$, the relation~\eqref{eqdiv3cc} is equivalent to  $q\equiv 0,2,4,5, 6\pmod{8}$.\ If the case $q\equiv 5\pmod{8}$ occurs, the same reasoning as in the case $n_X=1$,  $q\equiv 7\pmod{8}$, gives   a contradiction, hence $q_X$ is even.
\item When $n_X=4$, the relation~\eqref{eqdiv3cc} is equivalent to  $q\equiv 0,2,4,6\pmod{8}$ and $q_X$ is even.
 \end{itemize}

\medskip

\noindent{\bf{Assume  $ q_X(\lll,\mm)=2 $.}}\ The quadratic form $q_X$ is even, we have $c_X=\frac{15}4$ from~\eqref{cxa}, Lemma~\ref{le2} gives $m_X\in   \Z$, and~\eqref{mxroot} gives $m_X<5.8 $, so that $m_X \in\{1,2,3,4,5\}$.\  As above, we  deduce from~\eqref{prr3} that
 $$ \frac{1}{8\cdot 24}(2q+ 2m_X)^3+\Bigl(\frac{2}{m_X}-\frac{1}{8\cdot 24}4m_X^2\Bigr)(2q+ 2m_X) $$
 must be an integer for all values $2q$ taken by $q_X$, so that 
$$24 m_X\mid    m_X(q+m_X)^3+ (96- m_X^3 )(q+m_X).
$$
We reason as above to conclude that the integer~$q$ must be even,  so that all values taken by the quadratic form~$q_X$ are divisible by $4$.\ This is   impossible because   the gcd of all values taken  by $q_X$ is  $2$.\ So this case does not occur. 
 \end{proof}

 


\appendix

\section{Positivity of the Huybrechts--Riemann--Roch polynomial}\label{secapp}

\smallskip
\begin{center}by \textsc{Chen Jiang}\end{center}
\medskip

\newcommand\sigmabar{{\overline{\sigma}}}
\newcommand{\rounddown}[1]{\lfloor{#1}\rfloor}
\newcommand{\roundup}[1]{\lceil{#1}\rceil}
 

  Throughout this appendix, $X$ is a \hKm\ of dimension $2n$ and we fix a symplectic form $\sigma\in H^0(X, \Omega^2_X)$.\   The degree $n$  Huybrechts--Riemann--Roch polynomial  $P_{RR,X}(T)$ was defined in the introduction, and the polynomial $Q_{RR,X}(T)=P_{RR,X}(m_XT)$  in Section~\ref{sec:intro}.\ These polynomials were  proved in \cite[Theorem~1.1]{jia} to have positive coefficients.\ The purpose of this appendix is to prove a refinement of this result.\ For every nonnegative integer $k$, we define a degree $k$ monic polynomial with positive coefficients by
 $$Q_k(T)\coloneqq \sum_{j=0}^{k}\binom{k+j+1}{2j+1}T^j
 =T^k+2kT^{k-1}+\cdots + k+1
 .$$
  Our result is the following.
  
  \begin{prop}\label{propa1}
Let $X$ be a \hKm\ of dimension $2n>2$.\ There are nonnegative rational numbers $b_0, b_1, \dots, b_{\lfloor n/2\rfloor}$ 
such that
\begin{equation}\label{eq Q=sum Q}
Q_{RR, X}(T)=\sum_{i=0}^{\lfloor n/2\rfloor}b_iQ_{n-2i}(T).
\end{equation}
Moreover, $b_0=\int_X\td^{1/2}(X)>0$ and $b_1>0$. 
\end{prop}


 
  For any $\alpha\in H^2(X,\R)$, we have
\[
Q_{RR,X}(\lambda_X(\alpha))=\int_X \td(X) \exp(\alpha),
\]
where $\lambda_X$ is the quadratic form  on~$H^2(X,\R)$ discussed in Section~\ref{sec:intro}.\ Indeed, by~\eqref{hrrq}, this equality holds when $\alpha$ is the first Chern class of a line bundle on $X$.\ It then holds for each $\alpha\in H^2(X,\Z)$ because there is a deformation of $X$ on which $\alpha$ becomes the first Chern class of a line bundle.\ Finally, it holds for every $\alpha\in H^2(X,\R)$ since both sides are polynomial functions of $\alpha$.

Moreover, one has (\cite[Definition~17]{nw}, \cite[Definition~2.2]{jia})
\[ 
\lambda_X(\alpha)\coloneqq\begin{cases}\frac{24n \int_X \exp(\alpha)}{ \int_X c_{2}(X) \exp(\alpha)} & \text{if well-defined;}\\ 0 & \text{otherwise.}\end{cases} 
\]
For simplicity, we set $\lambda_\sigma\coloneqq\lambda_X(\sigma+\sigmabar)$. We know that $\lambda_\sigma>0$ (see \cite[Lemma~2.4(2)]{jia}). 

In \cite[Definition~4.1]{jia}, for any $0\leq k\leq n/2$,  
%$$
%\tp_{2k}\coloneqq\sum_{i=0}^k\frac{(n-2k+1)!\td^{1/2}_{2i}\wedge(\sigma\sigmabar)^{k-i}}{(-\lambda_\sigma)^{k-i}(k-i)!(n-k-i+1)!}\in H^{4k}(X).$$
%Note that it is of type $(2k,2k)$ with respect to the usual Hodge decomposition. 
we defined a class 
\[
\tp_{2k}\coloneqq\sum_{i=0}^k\frac{(n-2k+1)!\td^{1/2}_{2i}\wedge(\sigma\sigmabar)^{k-i}}{(-\lambda_\sigma)^{k-i}(k-i)!(n-k-i+1)!}\in H^{4k}(X,\R)\] 
which is of Hodge type $(2k,2k)$.\ 
%Recall that by Huybrechts' theorem \cite[Corollary~23.17]{GHJ}, we have constants ${\bf D}((\tp_{2i})^2)$ such that 
%$$\int_X (\tp_{2i})^2\alpha^{2n-4i}={\bf D}((\tp_{2i})^2)\cdot \lambda_X(\alpha)^{n-2i}$$
%for any $i$ and any $\alpha\in H^2(X)$.
One important fact is that, by \cite[Corollary~4.4]{jia}, 
$$\int_X\tp_{2k}^2(\sigma\sigmabar)^{n-2k}\geq 0.$$

\begin{lemm}\label{lem C independent}
The numbers
\[
C_k\coloneqq\frac{\int_X\tp_{2k}^2(\sigma\sigmabar)^{n-2k}}{\lambda_{\sigma}^{n-2k}}
\]
are deformation invariants of $X$.\ In particular, $C_k$ is independent of the choice of $\sigma$.
\end{lemm}

Here we remark that we cannot directly apply \cite[Corollary~23.17]{huy} as $\tp_{2k}$ might no longer be of type $(2k,2k)$ on deformations of $X$.

\begin{proof}
By  definition of $\tp_{2k}$, the number $C_k$ can be written as
\[C_k=\sum_{i=0}^k\sum_{j=0}^ka_{ij}\,\frac{\int_X\td^{1/2}_{2i}\td^{1/2}_{2j}(\sigma\sigmabar)^{n-i-j}}{\lambda_{\sigma}^{n-i-j}},\]
where the $a_{ij}$ are constants depending only on $n,k,i,j$.\ By \cite[Corollary~23.17]{huy} and \cite[Proposition~2.3]{jia},  \[\frac{\int_X\td^{1/2}_{2i}\td^{1/2}_{2j}(\sigma\sigmabar)^{n-i-j}}{\lambda_{\sigma}^{n-i-j}}=\frac{(n-i-j)!^2}{(2n-2i-2j)!}\,\frac{\int_X\td^{1/2}_{2i}\td^{1/2}_{2j}(\sigma+\sigmabar)^{2n-2i-2j}}{\lambda_{\sigma}^{n-i-j}}\]
 only depends  on $\td^{1/2}_{2i}\td^{1/2}_{2j}$, $c_2(X)$, and $c_X$, which implies that $C_k$ is a  deformation invariant of $X$. 
\end{proof}




\begin{proof}[Proof of Proposition~\ref{propa1}]
From \cite[Proof of Theorem~5.1]{jia}, for any $0\leq m\leq n$, we have \begin{align*}
 \int_X\td_{2m}(\sigma\sigmabar)^{n-m} 
= \sum_{i=0}^{\rounddown{m/2}} \frac{(n-m)!^2}{\lambda_\sigma^{m-2i} (n-2i)!^2}\binom{2n-2i-m+1}{m-2i}\int_X(\tp_{2i})^2(\sigma\sigmabar)^{n-2i}.
\end{align*}
In other words, 
\begin{align*}
 \int_X\td_{2m}(\sigma+\sigmabar)^{2n-2m} 
= \sum_{i=0}^{\rounddown{m/2}} \frac{(2n-2m)!}{\lambda_\sigma^{m-2i} (n-2i)!^2}\binom{2n-2i-m+1}{m-2i}\int_X(\tp_{2i})^2(\sigma\sigmabar)^{n-2i}.
\end{align*}

Thus we have the following equalities:
 \begin{align*}
\int_X \td(X) \exp(\sigma+\sigmabar)
 =&\sum_{m=0}^n\int_X  \frac{1}{(2n-2m)!}\td_{2m}(X) (\sigma+\sigmabar)^{2n-2m}\\
 =&\sum_{m=0}^n\sum_{i=0}^{\lfloor m/2\rfloor}\frac{1}{\lambda_\sigma^{m-2i}(n-2i)!^2}\binom{2n-2i-m+1}{m-2i}\int_X (\tp_{2i})^2(\sigma\sigmabar)^{n-2i}\\
 =&\sum_{m=0}^n\sum_{i=0}^{\lfloor m/2\rfloor}\frac{1}{(n-2i)!^2}\binom{2n-2i-m+1}{m-2i}C_i \lambda_\sigma^{n-m}\\
=&\sum_{i=0}^{\lfloor n/2\rfloor}\frac{C_i}{(n-2i)!^2} \sum_{m=2i}^n\binom{2n-2i-m+1}{m-2i}\lambda_\sigma^{n-m}\\
=&\sum_{i=0}^{\lfloor n/2\rfloor}\frac{C_i}{(n-2i)!^2} \sum_{m=0}^{n-2i}\binom{2n-4i-m+1}{m}\lambda_\sigma^{n-m-2i}\\
=&\sum_{i=0}^{\lfloor n/2\rfloor}\frac{C_i}{(n-2i)!^2} \,Q_{n-2i}(\lambda_\sigma).
\end{align*}
In other words,  
$$ Q_{RR, X}(\lambda_\sigma)=\sum_{i=0}^{\lfloor n/2\rfloor}\frac{C_i}{(n-2i)!^2} \,Q_{n-2i}(\lambda_\sigma).
$$
Here  $C_i\geq 0$ by \cite[Corollary~4.4]{jia}.\ 
By Lemma~\ref{lem C independent}, $C_i$ is independent of the choice of $\sigma$, so after replacing $\sigma$ by $t\sigma$ for any $t\in \C^\times$, we can get an equality of polynomials 
$$ Q_{RR, X}(T)=\sum_{i=0}^{\lfloor n/2\rfloor}\frac{C_i}{(n-2i)!^2}\, Q_{n-2i}(T),
$$
which gives  the desired equation \eqref{eq Q=sum Q}.
%So this observation gives another interpretation of  \cite[Corollary~5.2]{jia}, which says that the Riemann--Roch formula in terms of $\lambda$ can be written as an linear combination of $Q_{n-2i}$, $0\leq i\leq \lfloor n/2\rfloor$, with nonnegative coefficients.  

The last assertion is a consequence of \cite[Corollary~5.2]{jia}.
\end{proof}

 
 
 \begin{thebibliography}{DHMV}
 
 \bibitem[B]{bea} Beauville, A., Vari\'et\'es k\"ahl\'eriennes dont la premi\`ere classe de Chern est
nulle, {\em J. Differential Geom.} {\bf18} (1983),755--782. 
 
 
  \bibitem[BS]{beso}   Beckmann, T., Song, J.,
   Second Chern class and Fujiki constants of hyperk\"ahler manifolds, eprint {\tt arXiv:2201.07767.}
 
  \bibitem[DHMV]{dhmv} Debarre,  O., Huybrechts, D.,  Macr\`i, E.,   Voisin, C., Computing Riemann--Roch polynomials and classifying hyper-K\"ahler fourfolds, {\em J. Amer. Math. Soc.} {\bf37}  (2024), 151--185. 
  
   \bibitem[F]{fuj} Fujiki, A., 
On the de Rham cohomology group of a compact K\"ahler symplectic manifold, in {\em Algebraic geometry, Sendai, 1985,} 105--165, Adv. Stud. Pure Math. {\bf10}
North-Holland Publishing Co., Amsterdam, 1987.

%\bibitem[GHJ]{GHJ} Gross, M., Huybrechts,  D., Joyce, D., in {\it Calabi-Yau manifolds and related geometries}, Universitext, Springer-Verlag, Berlin, 2003. Lectures from the Summer School held in Nordfjordeid, June 2001.

  
 \bibitem[G]{gua}
Guan, D., On the Betti numbers of irreducible compact hyperk\"ahler manifolds of complex dimension four, {\em Math. Res. Lett.} {\bf8} (2001),   663--669.
 
 \bibitem[HS]{hs} Hitchin, N., Sawon, J., Curvature and characteristic numbers of hyper-K\"ahler manifolds, {\it Duke Math.~J.} {\bf 106} (2001), 599--615.

\bibitem[H1]{huyf} Huybrechts, D.,
Finiteness results for compact hyperk\"ahler manifolds, {\it J. reine angew. Math.} {\bf558} (2003), 15--22. 

\bibitem[H2]{huy} Huybrechts, D., Compact hyperk\"ahler manifolds, in {\em Calabi-Yau manifolds and related geometries (Nordfjordeid, 2001),} 161--225, Universitext, Springer, Berlin, 2003.
 
 \bibitem[J]{jia} Jiang, C., Positivity of Riemann--Roch polynomials and Todd classes of hyperk\"ahler manifolds, {\em J. Algebraic Geom.} {\bf32} (2023),  239--269.

 
 \bibitem[N]{nw} Nieper-Wisskirchen, M., Hirzebruch-Riemann-Roch formulae on irreducible symplectic K\"ahler manifolds, {\em J. Algebraic Geom.} {\bf12} (2003),  715--739.


 \end{thebibliography}

\end{document}


%%%%%%%%%%%%%%%%%%%%%

\begin{thebibliography}{Mat2}

%\bibitem[BHPV]{bpv} Barth, W., Hulek, K., Peters, Ch., Van de Ven, A., {\it Compact complex surfaces,} Second edition. Ergebnisse der Mathematik und ihrer Grenzgebiete {\bf4}, Springer-Verlag, Berlin, 2004.

\bibitem[BM1]{bm1}  Bayer, A., Macr\`i, E., Projectivity and birational geometry of Bridgeland moduli spaces, {\em J. Amer. Math. Soc.}  {\bf27} (2014),  707--752.

\bibitem[BM2]{bm}   Bayer, A., Macr\`i, E.,  MMP for moduli of sheaves on K3s via wall-crossing: nef and movable cones, Lagrangian fibrations, {\it Invent. Math.}  {\bf 198} (2014),  505--590.

\bibitem[B]{bea} Beauville, A.,
Counting rational curves on K3 surfaces, {\em
Duke Math. J.} {\bf97} (1999), 99--108. 
%
%\bibitem[B]{ben} Benoist, O.,  Espaces de modules d'intersections compl\`etes lisses, Th\`ese, 2012, available at {\tt https://www.math.ens.fr/$\sim$benoist/articles/These.pdf}
%
%\bibitem[BKPS]{bkps}  Borcherds, R., Katzarkov, L., Pantev, T., Shepherd-Barron, N.I., Families of K3 surfaces, {\it J. Algebraic Geom.}  {\bf7} (1998), 183--193. 

\bibitem[Bo]{bou}  Boucksom, S., Divisorial Zariski decompositions on compact complex manifolds, {\em Ann. Sci. \'Ecole
Norm. Sup.} {\bf37} (2004),  45--76.


%
%\bibitem[Br]{bra} Brakkee, E., Moduli spaces of twisted K3 surfaces and cubic fourfolds, {\it Math. Ann.} {\bf 377} (2020), 1453--1479.

\bibitem[Br]{bri:K3} Bridgeland, T., Stability conditions on K3 surfaces, {\it Duke Math. J.} {\bf 141} (2008), 241--291.

\bibitem[Bri]{bri}  Britze, M., Generalized Kummer varieties, Ph .D. thesis, University of K\"oln,
2002.

\bibitem[C]{cam} Campana, F., R\'eduction alg\'ebrique d'un morphisme faiblement K\"ahl\'erien propre et applications, {\it Math. Ann.} {\bf 256} (1981), 157--189.

 \bibitem[D]{kum} Debarre,  O.,  On the Euler characteristic of generalized Kummer varieties, {\em Amer. J. Math.} {\bf121} (1999), 577--586.

% \bibitem[D2]{debc} Debarre,  O.,   Cohomological characterizations of the complex projective space, eprint {\tt arXiv:1512.04321.}



 \bibitem[Dr1]{dru} Druel, S., Quelques remarques sur la d\'ecomposition de Zariski divisorielle sur les vari\'et\'es dont la
premi\`ere classe de Chern est nulle, {\em Math. Z.}  {\bf267} (2011),  413--423.

\bibitem[Dr2]{druel} Druel, S., A decomposition theorem for singular spaces with trivial canonical class of dimension at most five, {\it Invent. Math.}  {\bf211} (2018),  245--296.

\bibitem[EGL]{egl} Ellingsrud, G.,  G\"ottsche, L., Lehn, M.,
On the cobordism class of the Hilbert scheme of a surface, {\it 
J. Algebraic Geom.} {\bf10} (2001),  81--100. 

\bibitem[F]{fujiki} Fujiki, A., Relative algebraic reduction and relative Albanese map for a fiber space in $\cC$, {\it Publ. Res. Inst. Math. Sci.} {\bf 19} (1983), 207--236.

\bibitem[Fu]{fuj} Fujino, O., On Kawamata's theorem, in {\it Classification of algebraic varieties,} 305--315, EMS Ser. Congr. Rep., Eur. Math. Soc., Z\"urich, 2011.

%\bibitem[Fuk]{fuk} Fukuda, S., 

%
%\bibitem[GS1]{gs}  Garbagnati, A. Sarti, A., Projective models of K3 surfaces with an even set, {\it Adv. Geom.}  {\bf8} (2008), 413--440.
%
%\bibitem[GS2]{gs2}  Garbagnati, A. Sarti, A., Kummer surfaces and K3 surfaces with $(\Z/2\Z)^4$ symplectic action, {\it Rocky Mountain J. Math.}  {\bf46} (2016),  1141--1205.
%
%\bibitem[vdGK]{gk} van der Geer, G., Katsura, T., Note on tautological classes of moduli of K3 surfaces, {\it Mosc. Math. J.} {\bf5} (2005), 775--779, 972.
%
%\bibitem[GH]{gh} Griffiths, P., Harris, J., {\it Principles of algebraic geometry,} Pure and Applied Mathematics, Wiley-Interscience, New York, 1978.
% 
%% \bibitem[GrH]{grihul} Minimal Siegel modular threefolds, {\it Math. Proc. Cambridge Philos. Soc.} {\bf123} (1998), 461--485. 



%
%\bibitem[HM]{hamu} Hashimoto, K., Murabayashi, N., Shimura curves as intersections of Humbert surfaces and
%defining equations of QM-curves of genus two, {\it T\^ohoku Math. J.} {\bf47} (1995), 271--296.

\bibitem[HT]{hats} Hassett, B., Tschinkel, Y., Moving and ample cones of holomorphic symplectic fourfolds, {\it Geom.~ Funct.~ Anal.} {\bf 19} (2009), 1065--1080.

 
\bibitem[H1]{huyk3} Huybrechts, D., {\it Lectures on K3 surfaces,} Cambridge Studies in Advanced Mathematics {\bf158},  Cambridge University Press, 2016.

\bibitem[H2]{huy} Huybrechts, D., Compact hyperk\"ahler manifolds, in {\em Calabi-Yau manifolds and related geometries (Nordfjordeid, 2001),} 161--225, Universitext, Springer, Berlin, 2003.

\bibitem[H3]{huyf} Huybrechts, D.,
Finiteness results for compact hyperk\"ahler manifolds, {\it J. reine angew. Math.} {\bf558} (2003), 15--22. 

\bibitem[H4]{huyKahler} Huybrechts, D., ...

\bibitem[H5]{HuyHK} Huybrechts, D. {\it Compact hyperk\"ahler manifolds: Basic results,} Invent. Math. 135 (1999), 63--113.

\bibitem[HX]{hx} Huybrechts, D., Xu, C., 
Lagrangian fibrations of hyperk\"ahler fourfolds, eprint {\tt arXiv:1902.10440.}

\bibitem[J]{jia} Jiang, C., Positivity of Riemann--Roch polynomials and Todd classes of hyperk\"ahler manifolds, eprint {\tt  	arXiv:2008.04685.}
 
%
%\bibitem[HS]{HuSt} Huybrechts, D., Stellari, P., ...

\bibitem[K]{kaw} Kawamata, Y., Pluricanonical systems on minimal algebraic varieties, {\it Invent. Math.} {\bf79} (1985),  567--588.

\bibitem[KMM]{kmm}
Kawamata, Y.,  Matsuda, K.,  Matsuki, K.,
Introduction to the minimal model problem, in {\it Algebraic geometry, Sendai, 1985,} 283--360,
Adv. Stud. Pure Math. {\bf10}, North-Holland, Amsterdam, 1987.

%\bibitem[Ko]{kol} Koll\'ar, J., 

\bibitem[LSV]{lsv}    Laza, R.,  Sacc\`a, G.,  Voisin, C.,  A hyper-K\"ahler compactification
of the intermediate Jacobian fibration associated with a
cubic 4-fold, {\em Acta Math.} {\bf218} (2017), 55--135.

     \bibitem[L]{linthese}  Lin, H.-Y., Lagrangian constant cycle subvarieties in Lagrangian fibrations, {\it 
Int. Math. Res. Not. IMRN} 2020,   14--24.

% 
%\bibitem[L]{loo}  Looijenga, E., Compactifications defined by arrangements. II. Locally symmetric varieties of type IV,  {\it Duke Math. J.}  {\bf 119} (2003),  527--588.

 \bibitem[M]{mar}  Markman, E.,  Prime exceptional divisors on holomorphic symplectic varieties and monodromy-reflections, {\it Kyoto J. Math.} {\bf 53} (2013), 345--403.
 
     \bibitem[Ma]{martens} Martens, G., 
 
 \bibitem[Mat1]{mathigher} Matsushita, D., Higher direct images of dualizing sheaves of Lagrangian fibrations, {\it Amer. J. Math.} {\bf 127} (2005), 243--259.
 
 \bibitem[Mat2]{mat0}   Matsushita, D., On fibre space structures of a projective irreducible symplectic manifold,

 \bibitem[Mat3]{mat}   Matsushita, D., Addendum: ``On fibre space structures of a projective irreducible symplectic manifold,'' {\em Topology}  {\bf40}  (2001),  431--432.
 
   \bibitem[Mat4]{matlag}   Matsushita, D.,  On deformations of Lagrangian fibrations, in {\em K3 Surfaces and Their Moduli,} 237--243, Progr. Math. {\bf315}, Birkh\"auser/Springer, Cham, 2016.
  
 \bibitem[MO]{mo}   Mongardi, G.,  Onorati, C., Birational geometry of irreducible holomorphic symplectic tenfolds of O'Grady type, eprint {\tt  	arXiv:2008.04685.}

 
% 
%\bibitem[M]{muk} Mukai, S., Duality of polarized K3 surfaces, in {\it New trends in algebraic geometry (Warwick, 1996),} London Math. Soc. Lecture Note Ser. {\bf 264} (1999), Cambridge Univ. Press, 311--326.

 \bibitem[Mu]{mum}  Mumford, D. Pathologies. III, {\it Amer. J. Math.}. {\bf89} (1967), 94--104.

\bibitem[N]{nw} Nieper-Wisskirchen, M., Hirzebruch-Riemann-Roch formulae on irreducible symplectic K\"ahler manifolds, {\em J. Algebraic Geom.} {\bf12} (2003),  715--739.

%\bibitem[O]{ogu} Oguiso, K.,

\bibitem[O]{og1}	O'Grady,  K., Irreducible symplectic 4-folds numerically equivalent to $(K3)^{[2]}$,  {\it
Commun. Contemp. Math.}  {\bf10} (2008), 553--608. 

\bibitem[P]{peternell}  Peternell, T., Minimal varieties with trivial canonical classes, I, {\it Math. Z.}  {\bf217} (1994), 377--405.

\bibitem[R1]{rap} Rapagnetta, A., 
Topological invariants of O'Grady's six dimensional irreducible symplectic variety., {\em Math. Z.}  {\bf256} (2007),  1--34. 

\bibitem[R2]{rap2} Rapagnetta, A., On the Beauville form of the known irreducible symplectic varieties, {\em Math. Ann.} {\bf340} (2008), 77--95.

\bibitem[Ri]{ort} R\'ios Ortiz, A., Riemann-Roch Polynomials of the known Hyperk\"ahler Manifolds, eprint {\tt arXiv:2006.09307.}

\bibitem[SV]{sac} Sacc\`a, G., Birational geometry of the intermediate Jacobian fibration of a cubic fourfold, with an appendix by Claire Voisin, eprint {\tt arXiv:2002.01420.}

%\bibitem[S]{sal} Salamon, S. M. Spinors and cohomology, {\em Rend. Sem. Mat. Univ. Politec. Torino}  {\bf50} (1992),  393--410.

\bibitem[Sa]{saw} Sawon, J., A finiteness theorem for Lagrangian fibrations, {\em J. Algebraic Geom.} {\bf25} (2016),  431--459.

\bibitem[V1]{voi}  Voisin, C.,
Sur la stabilit\'e des sous-vari\'et\'es lagrangiennes des vari\'et\'es symplectiques holomorphes, in {\it Complex projective geometry (Trieste, 1989/Bergen, 1989),} 294--303,
London Math. Soc. Lecture Note Ser. {\bf179}, Cambridge Univ. Press, Cambridge, 1992.

 
\bibitem[V2]{voisinfib}  Voisin, C., On coverings and fibrations of hyper-K\"{a}hler manifolds, in preparation.

\bibitem[W]{wie} Wierzba, J.,  Contractions of symplectic varieties, {\it J. Algebraic Geom.}  {\bf12} (2003), 507--534.

\bibitem[WW]{wiwi} Wierzba, J., Wi\'sniewski, J. Small contractions and symplectic 4-folds, {\em Duke Math. J.} {\bf 120} (2003), 65--95.

\bibitem[Y]{yos} Yoshioka, K., Bridgeland's stability and the positive cone of the moduli spaces of stable objects on an abelian surface, in {\em Development of moduli theory, Kyoto 2013,} 473--537, Adv. Stud. Pure Math. {\bf69}, Math. Soc. Japan, Tokyo, 2016.

\end{thebibliography}

\end{document}

